% =====================================================================
\chapter{Integração}

No capítulo~2 vimos que o problema da área sob uma curva leva a um limite de somas de
retângulos. Neste capítulo esse limite ganha nome, \emph{integral definida}, e descobrimos o
fato central do Cálculo: integrar e derivar são operações inversas (Teorema Fundamental do
Cálculo). Com ele, calcular áreas, deslocamentos e variações totais se reduz a encontrar
antiderivadas, e a técnica de substituição amplia muito a lista de funções que sabemos
integrar.

% ---------------------------------------------------------------------
\section{Aproximando áreas}

\subsection{Notação de somatório}\index{somatorio@somatório}

Escrevemos $\displaystyle\sum_{i=1}^n a_i=a_1+a_2+\dots+a_n$. Valem
$\sum(a_i\pm b_i)=\sum a_i\pm\sum b_i$ e $\sum c\,a_i=c\sum a_i$, além das fórmulas
\[
\sum_{i=1}^n1=n,\qquad \sum_{i=1}^n i=\frac{n(n+1)}2,\qquad
\sum_{i=1}^n i^2=\frac{n(n+1)(2n+1)}6,\qquad \sum_{i=1}^n i^3=\frac{n^2(n+1)^2}4 .
\]

\subsection{Somas de Riemann}

Para aproximar a área sob o gráfico de $f\ge0$ em $[a,b]$, dividimos o intervalo em $n$
subintervalos de largura $\Delta x=\frac{b-a}n$, com extremos $x_i=a+i\,\Delta x$, e somamos
áreas de retângulos. A altura de cada retângulo é o valor de $f$ em algum ponto $x_i^*$ do
subintervalo:
\[
\sum_{i=1}^n f(x_i^*)\,\Delta x \qquad(\textbf{soma de Riemann}\index{soma de riemann@soma de Riemann}).
\]
Usando a extremidade esquerda ($x_{i-1}$), a direita ($x_i$) ou o ponto médio, obtemos as
aproximações $L_n$, $R_n$ e $M_n$.

\figtikz{\begin{tikzpicture}
\begin{groupplot}[group style={group size=3 by 1,horizontal sep=1.0cm},painel,width=5.3cm,height=4.6cm,
  xmin=-0.1,xmax=2.2,ymin=0,ymax=4.6,xtick={0,0.5,1,1.5,2},ytick={1,4},xticklabel style={font=\tiny}]
\nextgroupplot[title={$L_4=1{,}75$}]
\foreach \a in {0,0.5,1,1.5}{\edef\tmp{\noexpand\draw[fill=laranja!25,draw=laranja] (axis cs:\a,0) rectangle (axis cs:{\a+0.5},{\a*\a});}\tmp}
\addplot[domain=0:2.1]{x^2};
\nextgroupplot[title={$R_4=3{,}75$}]
\foreach \a in {0,0.5,1,1.5}{\edef\tmp{\noexpand\draw[fill=laranja!25,draw=laranja] (axis cs:\a,0) rectangle (axis cs:{\a+0.5},{(\a+0.5)^2});}\tmp}
\addplot[domain=0:2.1]{x^2};
\nextgroupplot[title={$M_4=2{,}625$}]
\foreach \a in {0,0.5,1,1.5}{\edef\tmp{\noexpand\draw[fill=laranja!25,draw=laranja] (axis cs:\a,0) rectangle (axis cs:{\a+0.5},{(\a+0.25)^2});}\tmp}
\addplot[domain=0:2.1]{x^2};
\end{groupplot}\end{tikzpicture}}{Somas de Riemann com $n=4$ para $f(x)=x^2$ em $[0,2]$:
extremidade esquerda, direita e ponto médio. A área exata é $\tfrac83\approx2{,}667$.}

\begin{exemplo}{5.1}{Aproximações e o limite}
Para $f(x)=x^2$ em $[0,2]$ com $n=4$ ($\Delta x=0{,}5$):
$L_4=0{,}5(0+0{,}25+1+2{,}25)=1{,}75$ e $R_4=0{,}5(0{,}25+1+2{,}25+4)=3{,}75$. Como $f$ é
crescente, $L_n$ subestima e $R_n$ superestima a área.

Para $n$ qualquer, $x_i=\frac{2i}n$ e
\[
R_n=\sum_{i=1}^n\Big(\frac{2i}n\Big)^2\frac2n=\frac8{n^3}\sum_{i=1}^ni^2
=\frac8{n^3}\cdot\frac{n(n+1)(2n+1)}6\ \xrightarrow{n\to\infty}\ \frac{16}{6}=\frac83 .
\]
\end{exemplo}

% ---------------------------------------------------------------------
\section{A integral definida}

\begin{definicao}
Se $f$ está definida em $[a,b]$ e o limite
\[
\int_a^b f(x)\,dx=\lim_{n\to\infty}\sum_{i=1}^n f(x_i^*)\,\Delta x
\]
existe e é o mesmo para quaisquer escolhas dos pontos $x_i^*$, dizemos que $f$ é
\textbf{integrável}\index{funcao@função!integravel@integrável} em $[a,b]$ e o limite é a \textbf{integral definida}\index{integral!definida} de $f$ de $a$ a $b$.
Os números $a$ e $b$ são os \emph{limites de integração}, e $f$ é o \emph{integrando}.
\end{definicao}

Toda função contínua em $[a,b]$ é integrável (e também as que têm um número finito de
descontinuidades de salto). A variável de integração é ``muda'':
$\int_a^bf(x)\,dx=\int_a^bf(t)\,dt$.

\textbf{Área com sinal.}\index{area com sinal@área com sinal} Se $f$ muda de sinal, os retângulos abaixo do eixo contribuem
negativamente: $\int_a^bf(x)\,dx$ é a \emph{área líquida} (área acima do eixo menos área
abaixo). A \emph{área total} é $\int_a^b|f(x)|\,dx$.

\figtikz{\begin{tikzpicture}
\begin{axis}[painel,width=9cm,height=4.4cm,xmin=-0.2,xmax=6.6,ymin=-1.3,ymax=1.3,
  xtick={3.1416,6.2832},xticklabels={$\pi$,$2\pi$},ytick={-1,1}]
\addplot[fill=azul!25,draw=none,domain=0:3.1416] {sin(deg(x))} \closedcycle;
\addplot[fill=laranja!30,draw=none,domain=3.1416:4.7124] {sin(deg(x))} \closedcycle;
\addplot[domain=-0.2:6.6]{sin(deg(x))};
\node[azul,font=\small] at (axis cs:1.57,0.4) {$+2$}; \node[laranja,font=\small] at (axis cs:3.95,-0.4) {$-1$};
\end{axis}\end{tikzpicture}}{$\int_0^{3\pi/2}\sen x\,dx=2-1=1$ (área líquida), enquanto a área
total é $2+1=3$.}

\begin{teorema}[Propriedades da integral definida]{5.1}
Para $f$ e $g$ integráveis e $k$ constante:
\begin{enumerate}[label=(\roman*)]
\item $\int_a^a f=0$ \ e \ $\int_b^a f=-\int_a^b f$;
\item $\int_a^b(f\pm g)=\int_a^bf\pm\int_a^bg$ \ e \ $\int_a^bkf=k\int_a^bf$;
\item $\int_a^bf=\int_a^cf+\int_c^bf$;
\item se $f\le g$ em $[a,b]$, então $\int_a^bf\le\int_a^bg$; se $m\le f\le M$, então
  $m(b-a)\le\int_a^bf\le M(b-a)$.
\end{enumerate}
\end{teorema}

\textbf{Valor médio.} O \textbf{valor médio}\index{valor medio de uma funcao@valor médio de uma função} de $f$ em $[a,b]$ é
$f_{\text{méd}}=\dfrac1{b-a}\displaystyle\int_a^bf(x)\,dx$: a altura do retângulo de base
$[a,b]$ com a mesma área líquida que a região sob $f$.

% ---------------------------------------------------------------------
\section{O Teorema Fundamental do Cálculo}

\begin{teorema}[Teorema do Valor Médio para integrais]{5.2}\index{teorema@Teorema!do valor medio para integrais@do Valor Médio para integrais}
Se $f$ é contínua em $[a,b]$, existe $c\in[a,b]$ tal que $f(c)=\dfrac1{b-a}\displaystyle\int_a^bf(x)\,dx$.
\end{teorema}

Considere a \emph{função área} $F(x)=\int_a^xf(t)\,dt$: a área acumulada de $a$ até $x$. Para
$h$ pequeno, $F(x+h)-F(x)=\int_x^{x+h}f(t)\,dt\approx f(x)\,h$ (uma fatia de largura $h$ e
altura $\approx f(x)$). Isso sugere que $F'(x)=f(x)$.

\figtikz{\begin{tikzpicture}
\begin{axis}[painel,width=8cm,height=4.8cm,xmin=0,xmax=5,ymin=0,ymax=3.6,xtick={0.5,2.6,3.0},xticklabels={$a$,$x$,$x\!+\!h$},ytick=\empty]
\addplot[fill=azul!20,draw=none,domain=0.5:2.6] {1+0.6*x-0.08*x^2+0.3*sin(deg(2*x))} \closedcycle;
\addplot[fill=laranja!40,draw=none,domain=2.6:3.0] {1+0.6*x-0.08*x^2+0.3*sin(deg(2*x))} \closedcycle;
\addplot[domain=0:5]{1+0.6*x-0.08*x^2+0.3*sin(deg(2*x))};
\node[azul,font=\small] at (axis cs:1.55,0.8) {$F(x)$};
\end{axis}\end{tikzpicture}}{A função área $F(x)=\int_a^xf(t)\,dt$. Ao passar de $x$ para
$x+h$, a área cresce cerca de $f(x)\cdot h$ (faixa verde).}

\begin{teorema}[Teorema Fundamental do Cálculo, parte 1]{5.3}\index{teorema@Teorema!fundamental do calculo@Fundamental do Cálculo}
Se $f$ é contínua em $[a,b]$, então $F(x)=\int_a^xf(t)\,dt$ é derivável em $(a,b)$ e
\[
F'(x)=\frac{d}{dx}\int_a^xf(t)\,dt=f(x).
\]
\end{teorema}
\emph{Prova.} Pelo TVM para integrais, $\frac{F(x+h)-F(x)}h=\frac1h\int_x^{x+h}f(t)\,dt=f(c_h)$
para algum $c_h$ entre $x$ e $x+h$. Quando $h\to0$, $c_h\to x$ e, por continuidade,
$f(c_h)\to f(x)$. $\square$

Combinada com a regra da cadeia:
$\dfrac{d}{dx}\displaystyle\int_a^{g(x)}f(t)\,dt=f\big(g(x)\big)\,g'(x)$. Por exemplo,
$\dfrac{d}{dx}\displaystyle\int_1^{x^2}\cos t\,dt=2x\cos(x^2)$.

\begin{teorema}[Teorema Fundamental do Cálculo, parte 2]{5.4}
Se $f$ é contínua em $[a,b]$ e $F$ é qualquer antiderivada de $f$, então
\[
\int_a^bf(x)\,dx=F(b)-F(a)=\Big[F(x)\Big]_a^b .
\]
\end{teorema}
\emph{Prova.} $G(x)=\int_a^xf(t)\,dt$ também é antiderivada de $f$ (parte 1), logo
$G=F+C$. Então $\int_a^bf=G(b)-G(a)=F(b)-F(a)$. $\square$

\begin{exemplo}{5.2}{Usando o TFC}
(a) $\int_0^2x^2\,dx=\Big[\frac{x^3}3\Big]_0^2=\frac83$, confirmando o Exemplo~5.1 sem limites.

(b) $\int_0^\pi\sen x\,dx=\big[-\cos x\big]_0^\pi=1+1=2$.

(c) $\displaystyle\int_1^4\Big(\sqrt x+\frac1{x^2}\Big)dx=\Big[\tfrac23x^{3/2}-\tfrac1x\Big]_1^4
=\Big(\tfrac{16}3-\tfrac14\Big)-\Big(\tfrac23-1\Big)=\tfrac{65}{12} .$
\end{exemplo}

% ---------------------------------------------------------------------
\section{Fórmulas de integração e o Teorema da Variação Total}

Lendo a parte 2 do TFC com $F'$ no lugar de $f$: a integral de uma taxa de variação é a
variação total.

\begin{teorema}[Teorema da Variação Total]{5.5}\index{teorema@Teorema!da variacao total@da Variação Total}
$F(b)=F(a)+\displaystyle\int_a^bF'(x)\,dx$, ou seja, $\displaystyle\int_a^bF'(x)\,dx=F(b)-F(a)$.
\end{teorema}

Se $v(t)$ é a velocidade de um objeto, $\int_a^bv(t)\,dt$ é o \textbf{deslocamento}\index{deslocamento} (variação
de posição), enquanto $\int_a^b|v(t)|\,dt$ é a \textbf{distância percorrida}\index{distancia percorrida@distância percorrida}. Se $r(t)$ é a
vazão de água num tanque, $\int_a^br(t)\,dt$ é a variação do volume.

\begin{exemplo}{5.3}{Deslocamento \texorpdfstring{$\times$}{x} distância}
Uma partícula tem velocidade $v(t)=t^2-4t+3=(t-1)(t-3)$ m/s, $0\le t\le4$.
\solucao
Deslocamento: $\int_0^4v\,dt=\big[\frac{t^3}3-2t^2+3t\big]_0^4=\frac43$ m.
Para a distância, separamos onde $v$ muda de sinal ($t=1$ e $t=3$):
$\int_0^1v=\frac43$, $\int_1^3v=-\frac43$, $\int_3^4v=\frac43$. Distância
$=\frac43+\frac43+\frac43=4$ m.
\end{exemplo}

\textbf{Funções pares e ímpares.}\index{simetria em integrais} Se $f$ é par, $\int_{-a}^af=2\int_0^af$; se $f$ é ímpar,
$\int_{-a}^af=0$ (as áreas se cancelam). Por exemplo, $\int_{-2}^2x^3\cos x\,dx=0$ sem nenhuma
conta.

% ---------------------------------------------------------------------
\section{Substituição}

A regra da cadeia diz que $\frac{d}{dx}F(g(x))=F'(g(x))\,g'(x)$. Lida ao contrário:
\[
\int f\big(g(x)\big)\,g'(x)\,dx=\int f(u)\,du,\qquad u=g(x),\ du=g'(x)\,dx .
\]

\begin{estrategia}{integração por substituição}\index{substituicao@substituição}
\begin{enumerate}[label=\arabic*.]
\item Escolha $u=g(x)$, em geral a ``função de dentro'', cuja derivada aparece (a menos de
  uma constante) no integrando.
\item Calcule $du=g'(x)\,dx$ e reescreva \emph{toda} a integral em termos de $u$.
\item Integre em $u$ e volte para $x$. Em integrais definidas, troque também os limites:
  $x=a\mapsto u=g(a)$ e $x=b\mapsto u=g(b)$.
\end{enumerate}
\end{estrategia}

\begin{exemplo}{5.4}{Substituições}
(a) $\int x^2(x^3+5)^9\,dx$: com $u=x^3+5$, $du=3x^2dx$, fica
$\frac13\int u^9\,du=\frac{u^{10}}{30}+C=\frac{(x^3+5)^{10}}{30}+C$.

(b) $\int\sen(5x)\,dx=-\frac15\cos(5x)+C$ \quad ($u=5x$).

(c) $\int x\sqrt{x-1}\,dx$: com $u=x-1$, temos $x=u+1$ e
$\int(u+1)u^{1/2}\,du=\frac25u^{5/2}+\frac23u^{3/2}+C=\frac25(x-1)^{5/2}+\frac23(x-1)^{3/2}+C$.

(d) $\int_0^{\pi/2}\sen^2x\cos x\,dx$: com $u=\sen x$, os limites viram $0$ e $1$:
$\int_0^1u^2\,du=\frac13$.
\end{exemplo}

% ---------------------------------------------------------------------
\section{Integrais com exponenciais e logaritmos}

Das derivadas do capítulo~3:
\[
\int e^u\,du=e^u+C,\qquad \int a^u\,du=\frac{a^u}{\ln a}+C,\qquad \int\frac{du}u=\ln|u|+C .
\]

\begin{exemplo}{5.5}{Exponenciais e logaritmos}
(a) $\int xe^{-x^2}dx=-\frac12\int e^u\,du=-\frac12e^{-x^2}+C$ \quad ($u=-x^2$).

(b) $\int\frac{3x^2}{x^3+1}\,dx=\ln|x^3+1|+C$: o numerador é a derivada do denominador.

(c) $\int\tg x\,dx=\int\frac{\sen x}{\cos x}\,dx=-\ln|\cos x|+C=\ln|\sec x|+C$ \quad
($u=\cos x$).

(d) $\int_1^e\frac{\ln x}x\,dx=\int_0^1u\,du=\frac12$ \quad ($u=\ln x$).
\end{exemplo}

% ---------------------------------------------------------------------
\section{Integrais que resultam em inversas trigonométricas}

Das derivadas das inversas trigonométricas, para $a>0$:
\[
\int\frac{du}{\sqrt{a^2-u^2}}=\arcsen\frac ua+C,\qquad
\int\frac{du}{a^2+u^2}=\frac1a\arctg\frac ua+C,\qquad
\int\frac{du}{u\sqrt{u^2-a^2}}=\frac1a\arcsec\frac{|u|}a+C .
\]

\begin{exemplo}{5.6}{Inversas trigonométricas}
(a) $\displaystyle\int\frac{dx}{\sqrt{4-9x^2}}$: com $u=3x$, $\frac13\int\frac{du}{\sqrt{4-u^2}}
=\frac13\arcsen\frac{3x}2+C$.

(b) $\displaystyle\int_0^1\frac{dx}{1+x^2}=\big[\arctg x\big]_0^1=\frac\pi4$. Daí,
$\pi=\int_0^1\frac{4}{1+x^2}\,dx$.

(c) $\displaystyle\int\frac{dx}{x^2+4x+13}$: completando o quadrado, $x^2+4x+13=(x+2)^2+9$, e
com $u=x+2$: $\frac13\arctg\frac{x+2}3+C$.
\end{exemplo}

% ---------------------------------------------------------------------
\section*{Resumo do capítulo}\markright{Resumo do capítulo}
\addcontentsline{toc}{section}{Resumo do capítulo}
\begin{itemize}
\item Somas de Riemann $\sum f(x_i^*)\Delta x$ aproximam a área; $\int_a^bf$ é o seu limite e
  mede a área com sinal.
\item Propriedades: linearidade, aditividade em intervalos, inversão de limites, comparação;
  valor médio $\frac1{b-a}\int_a^bf$.
\item TFC 1: $\frac d{dx}\int_a^xf(t)\,dt=f(x)$ (com a cadeia, $f(g(x))g'(x)$).
  TFC 2: $\int_a^bf=F(b)-F(a)$.
\item Variação total: $\int_a^bF'=F(b)-F(a)$. Deslocamento $=\int v$; distância
  $=\int|v|$.
\item Simetria: $\int_{-a}^a$(ímpar) $=0$; $\int_{-a}^a$(par) $=2\int_0^a$.
\item Substituição $u=g(x)$, $du=g'(x)dx$; nas definidas, troque os limites.
\item $\int e^u=e^u$, $\int\frac{du}u=\ln|u|$, $\int\frac{du}{\sqrt{a^2-u^2}}=\arcsen\frac ua$,
  $\int\frac{du}{a^2+u^2}=\frac1a\arctg\frac ua$.
\end{itemize}

% Gerado por ferramentas/gerar.py; não edite à mão.
\section*{Exercícios}\markright{Exercícios}
\addcontentsline{toc}{section}{Exercícios}
Os pontos antes de cada número indicam a dificuldade: \nivel{1} fácil, \nivel{2} médio, \nivel{3} difícil e \nivel{4} desafio. Os exercícios marcados com \estrela{} têm solução completa, com animação do gráfico, no repositório do livro: \url{https://github.com/AleffAzevedo/calculo-1-pt/blob/main/exercicios/cap5/README.md}. As respostas estão no fim do capítulo.

\begin{itemize}[leftmargin=5.2em,labelsep=0.6em,itemsep=4pt]
\item[\exnum{1}{5.1}] Calcule $\displaystyle\sum_{i=1}^{20}(3i-2)$ e $\displaystyle\sum_{i=1}^{10}i^2$ usando as fórmulas de somatório.
\item[\exnum{1}{5.2}] Calcule $L_4$ e $R_4$ para $f(x)=x^2+1$ em $[0,2]$. Qual subestima a área? Por quê?
\item[\exnum{1}{5.3}] Calcule a soma do ponto médio $M_4$ para $f(x)=\dfrac1x$ em $[1,3]$.
\item[\exnum{1}{5.4}] Interprete $\int_0^3\sqrt{9-x^2}\,dx$ como área e calcule-a sem antiderivadas.
\item[\exnum{1}{5.5}] Use áreas de figuras geométricas para calcular $\int_{-2}^{4}|x-1|\,dx$.
\item[\exnum{1}{5.6}] Sabendo que $\int_0^5f=7$, $\int_0^3f=4$ e $\int_0^5g=-2$, calcule $\int_3^5f$, $\int_5^0(2f-3g)$ e $\int_3^3g$.
\item[\exnum{1}{5.7}] Calcule o valor médio de $f(x)=x^2$ em $[0,3]$ e encontre $c$ com $f(c)=f_{\text{méd}}$.
\item[\exnum{1}{5.8}] Calcule $\dfrac{d}{dx}\displaystyle\int_2^x\frac{\sin t}{t}\,dt$ e $\dfrac{d}{dx}\displaystyle\int_0^{x^3}\sqrt{1+t^2}\,dt$.
\item[\exnum{1}{5.9}] Calcule $\displaystyle\int_1^2\Big(3x^2-\frac2{x^3}\Big)dx$.
\item[\exnum{1}{5.10}] Calcule $\displaystyle\int_0^{\pi/3}\sec^2x\,dx$ e $\displaystyle\int_0^{\pi}(\cos x+2)\,dx$.
\item[\exnum{1}{5.11}] Calcule $\displaystyle\int_1^4\frac{x-1}{\sqrt x}\,dx$.
\item[\exnum{1}{5.12}] Calcule $\displaystyle\int_0^1(e^x-2^x)\,dx$.
\item[\exnum{1}{5.13}] Calcule $\displaystyle\int(2x+1)^7\,dx$ e $\displaystyle\int\cos(3x-1)\,dx$.
\item[\exnum{1}{5.14}] Calcule $\displaystyle\int x\sqrt{x^2+1}\,dx$.
\item[\exnum{1}{5.15}] Calcule $\displaystyle\int\frac{x}{x^2+4}\,dx$ e $\displaystyle\int\frac{dx}{2x+5}$.
\item[\exnum{1}{5.16}] Calcule $\displaystyle\int e^{5x}\,dx$ e $\displaystyle\int x^2e^{x^3}\,dx$.
\item[\exnum{1}{5.17}] Calcule $\displaystyle\int\frac{dx}{\sqrt{9-x^2}}$ e $\displaystyle\int\frac{dx}{16+x^2}$.
\item[\exnum{1}{5.18}] Calcule $\displaystyle\int_0^{\pi/2}\sin x\cos^3x\,dx$.
\item[\exnum{1}{5.19}] Calcule $\displaystyle\int_{-1}^{1}(x^5-3x^3+x^2)\,dx$ usando simetria.
\item[\exnum{1}{5.20}] Uma partícula tem $v(t)=2t-4$ m/s em $[0,5]$. Calcule o deslocamento e a distância percorrida.
\item[\exnum{2}{5.21}] \estrela\ \textbf{Somas de Riemann convergindo.} Para $f(x)=4-x^2$ em $[0,2]$, calcule $L_4$ e $R_4$, escreva $R_n$ em função de $n$ e
encontre a área exata como $\lim_{n\to\infty}R_n$.
\item[\exnum{2}{5.22}] \estrela\ \textbf{Integral pela definição.} Use a definição (limite de somas de Riemann com extremidades direitas) para calcular
$\int_0^1x^3\,dx$.
\item[\exnum{2}{5.23}] Escreva $\displaystyle\lim_{n\to\infty}\sum_{i=1}^n\Big(1+\frac{2i}n\Big)^2\frac2n$ como integral definida e calcule.
\item[\exnum{2}{5.24}] Mostre, sem calcular, que $2\le\displaystyle\int_0^2\sqrt{1+x^3}\,dx\le6$.
\item[\exnum{2}{5.25}] \estrela\ \textbf{A função área.} Seja $g(x)=\displaystyle\int_0^x(t-1)(t-3)\,dt$ para $0\le x\le4$. Use o TFC para encontrar
$g'(x)$, os intervalos em que $g$ cresce e decresce, e os extremos locais de $g$. Calcule
$g(1)$, $g(3)$ e $g(4)$.
\item[\exnum{2}{5.26}] \estrela\ \textbf{TFC com limites variáveis.} Calcule $\dfrac{d}{dx}\displaystyle\int_x^{x^2}e^{t^2}\,dt$.
\item[\exnum{2}{5.27}] Calcule $\dfrac{d}{dx}\displaystyle\int_{\cos x}^{1}\frac{dt}{1+t^2}$.
\item[\exnum{2}{5.28}] Calcule $\displaystyle\int_0^3|x^2-4|\,dx$.
\item[\exnum{2}{5.29}] \estrela\ \textbf{Deslocamento e distância.} Uma partícula tem velocidade $v(t)=3t^2-12t+9$ m/s para $0\le t\le4$. Calcule o
deslocamento e a distância percorrida. Se $s(0)=2$, onde ela está em $t=4$?
\item[\exnum{2}{5.30}] \estrela\ \textbf{Valor médio de uma função.} Calcule o valor médio de $f(x)=\sin x$ em $[0,\pi]$ e encontre todos os $c$ garantidos pelo
TVM para integrais.
\item[\exnum{2}{5.31}] Calcule $\displaystyle\int x\sqrt{x+2}\,dx$.
\item[\exnum{2}{5.32}] Calcule $\displaystyle\int\frac{\cos x}{\sin^2x}\,dx$ e $\displaystyle\int\sec^2x\,\tan^3x\,dx$.
\item[\exnum{2}{5.33}] \estrela\ \textbf{Substituição em integral definida.} Calcule $\displaystyle\int_0^2x\sqrt{4-x^2}\,dx$ trocando os limites de integração, e interprete
a substituição geometricamente.
\item[\exnum{2}{5.34}] Calcule $\displaystyle\int_0^{1}\frac{e^x}{1+e^x}\,dx$.
\item[\exnum{2}{5.35}] Calcule $\displaystyle\int\frac{(\ln x)^2}{x}\,dx$ e $\displaystyle\int\frac{dx}{x\ln x}$.
\item[\exnum{2}{5.36}] Calcule $\displaystyle\int\frac{e^{x}}{1+e^{2x}}\,dx$.
\item[\exnum{2}{5.37}] Calcule $\displaystyle\int_0^{1/2}\frac{dx}{\sqrt{1-x^2}}$ e $\displaystyle\int_{0}^{\sqrt3}\frac{dx}{1+x^2}$.
\item[\exnum{2}{5.38}] Calcule $\displaystyle\int\frac{dx}{x^2-6x+13}$ completando o quadrado.
\item[\exnum{2}{5.39}] \estrela\ \textbf{Simetria salva contas.} Calcule $\displaystyle\int_{-\pi}^{\pi}\big(x^3\cos x+2\cos x+x\big)\,dx$ usando simetria.
\item[\exnum{2}{5.40}] Calcule $\displaystyle\int\frac{x+3}{x^2+1}\,dx$ (separe em duas integrais).
\item[\exnum{3}{5.41}] \estrela\ \textbf{Calculando $\pi$ com uma integral.} Mostre que $\displaystyle\int_0^1\frac{4}{1+x^2}\,dx=\pi$ e use a regra do ponto médio com $n=4$ para
estimar $\pi$.
\item[\exnum{3}{5.42}] \estrela\ \textbf{Vazão e volume.} Água entra num tanque à taxa $r(t)=200-4t$ L/min, $0\le t\le50$, que inicialmente tem
1000 L. Quanto há no tanque em $t=50$? Em que instante o volume cresce mais rapidamente?
Esboce $V(t)$.
\item[\exnum{3}{5.43}] \estrela\ \textbf{Integral da tangente e da secante.} Mostre que $\int\tan x\,dx=\ln|\sec x|+C$ e que $\int\sec x\,dx=\ln|\sec x+\tan x|+C$.
Calcule $\int_0^{\pi/4}\sec x\,dx$.
\item[\exnum{3}{5.44}] Use somas de Riemann (extremidades direitas) para calcular $\int_1^3(2x+1)\,dx$ pela definição.
\item[\exnum{3}{5.45}] Encontre $f$ contínua e a constante $a$ tais que $\displaystyle6+\int_a^x\frac{f(t)}{t^2}\,dt=2\sqrt x$ para $x>0$.
\item[\exnum{3}{5.46}] Encontre os extremos locais e os pontos de inflexão de $F(x)=\displaystyle\int_0^x\frac{t-2}{1+t^2}\,dt$, $x\ge0$.
\item[\exnum{3}{5.47}] Calcule $\displaystyle\lim_{x\to0}\frac1{x^3}\int_0^x\sin(t^2)\,dt$.
\item[\exnum{3}{5.48}] Calcule $\displaystyle\int_0^{\pi}|\cos x|\,dx$ e $\displaystyle\int_0^{2\pi}|\sin x-\cos x|\,dx$.
\item[\exnum{3}{5.49}] A taxa de vendas de um produto é $r(t)=100\,t\,e^{-t^2/10}$ unidades/semana. Quantas unidades são vendidas nas 4 primeiras semanas? E no longo prazo?
\item[\exnum{3}{5.50}] Calcule $\displaystyle\int\frac{x^3}{\sqrt{1-x^2}}\,dx$ (dica: $u=1-x^2$).
\item[\exnum{3}{5.51}] Calcule $\displaystyle\int\frac{dx}{\sqrt{e^{2x}-1}}$ (dica: $u=e^x$).
\item[\exnum{3}{5.52}] Calcule $\displaystyle\int\frac{\sin x}{1+\cos^2x}\,dx$ e $\displaystyle\int\frac{dx}{x\sqrt{x^2-4}}$.
\item[\exnum{3}{5.53}] Calcule $\displaystyle\int\frac{dx}{\sqrt{-x^2+4x}}$ completando o quadrado.
\item[\exnum{3}{5.54}] Mostre que, se $f$ é contínua e periódica de período $T$, então $\int_a^{a+T}f$ não depende de $a$.
\item[\exnum{3}{5.55}] Um carro freia com $a(t)=-8$ m/s² a partir de $30$ m/s. Use integrais para achar a distância até parar.
\item[\exnum{3}{5.56}] O valor eficaz (RMS) de $f$ em $[a,b]$ é $\sqrt{\tfrac1{b-a}\int_a^bf^2}$. Calcule o RMS de $f(t)=\sin t$ em $[0,2\pi]$. (Dica: $\sin^2t=\tfrac{1-\cos2t}2$.)
\item[\exnum{3}{5.57}] Seja $g(x)=\int_0^x f(t)\,dt$, com $f$ crescente e $f(0)=0$. Mostre que $g$ é côncava para cima e que $g(x)\ge0$ para $x\ge0$.
\item[\exnum{3}{5.58}] Calcule $\displaystyle\lim_{n\to\infty}\frac1n\sum_{i=1}^n\sqrt{\frac in}$ interpretando como integral.
\item[\exnum{3}{5.59}] Calcule $\displaystyle\lim_{n\to\infty}\sum_{i=1}^n\frac{n}{n^2+i^2}$.
\item[\exnum{3}{5.60}] Calcule $\displaystyle\int_0^{\ln2}e^x\sqrt{e^x-1}\,dx$ e $\displaystyle\int_1^{e^2}\frac{dx}{x\sqrt{\ln x}}$.
\item[\exnum{4}{5.61}] Prove que $\displaystyle\int_0^{\pi/2}\frac{\sin^nx}{\sin^nx+\cos^nx}\,dx=\frac\pi4$ para todo $n>0$. (Dica: $u=\tfrac\pi2-x$.)
\item[\exnum{4}{5.62}] Mostre que $\ln x=\int_1^x\frac{dt}{t}$ satisfaz $\ln(ab)=\ln a+\ln b$ usando apenas propriedades da integral e substituição.
\item[\exnum{4}{5.63}] Calcule $\displaystyle\int_0^1\frac{\ln(1+x)}{1+x^2}\,dx$. (Dica: $x=\tan\theta$ e depois $\theta\mapsto\tfrac\pi4-\theta$.)
\item[\exnum{4}{5.64}] Encontre todas as funções contínuas $f$ com $\displaystyle\int_0^xf(t)\,dt=x^2+\int_x^1t^2f(t)\,dt$ para todo $x$.
\item[\exnum{4}{5.65}] Prove a desigualdade $\ln(1+n)<1+\tfrac12+\dots+\tfrac1n<1+\ln n$ comparando somas com $\int\frac{dx}x$.
\end{itemize}
\subsection*{Respostas}\markright{Respostas}
{\small\begin{multicols}{2}\raggedright
\begin{itemize}[leftmargin=2.7em,labelsep=0.5em,itemsep=1pt]
\item[\textbf{5.1}] $590$ e $385$.
\item[\textbf{5.2}] $L_4=3{,}75$; $R_4=5{,}75$; $L_4$ subestima, pois $f$ é crescente.
\item[\textbf{5.3}] $M_4=\tfrac{2}{5}+\tfrac27+\tfrac29+\tfrac2{11}\approx1{,}0898$ (exato: $\ln3\approx1{,}0986$).
\item[\textbf{5.4}] Um quarto do círculo de raio 3: $\tfrac{9\pi}4$.
\item[\textbf{5.5}] $\tfrac92+\tfrac92=9$.
\item[\textbf{5.6}] $3$; $-20$; $0$.
\item[\textbf{5.7}] $3$; $c=\sqrt3$.
\item[\textbf{5.8}] $\tfrac{\sin x}x$; $3x^2\sqrt{1+x^6}$.
\item[\textbf{5.9}] $\tfrac{25}4$.
\item[\textbf{5.10}] $\sqrt3$; $2\pi$.
\item[\textbf{5.11}] $\tfrac83$.
\item[\textbf{5.12}] $e-1-\tfrac1{\ln2}\approx0{,}276$.
\item[\textbf{5.13}] $\tfrac{(2x+1)^8}{16}+C$; $\tfrac13\sin(3x-1)+C$.
\item[\textbf{5.14}] $\tfrac13(x^2+1)^{3/2}+C$.
\item[\textbf{5.15}] $\tfrac12\ln(x^2+4)+C$; $\tfrac12\ln|2x+5|+C$.
\item[\textbf{5.16}] $\tfrac15e^{5x}+C$; $\tfrac13e^{x^3}+C$.
\item[\textbf{5.17}] $\sin^{-1}\tfrac x3+C$; $\tfrac14\tan^{-1}\tfrac x4+C$.
\item[\textbf{5.18}] $\tfrac14$.
\item[\textbf{5.19}] $\tfrac23$.
\item[\textbf{5.20}] $5$ m e $13$ m.
\item[\textbf{5.21}] $L_4=6{,}25$; $R_4=4{,}25$; $R_n=8-\tfrac{4(n+1)(2n+1)}{3n^2}$; área $\tfrac{16}3$.
\item[\textbf{5.22}] $\tfrac14$.
\item[\textbf{5.23}] $\int_1^3x^2\,dx=\tfrac{26}3$.
\item[\textbf{5.24}] $1\le\sqrt{1+x^3}\le3$ em $[0,2]$.
\item[\textbf{5.25}] $g'(x)=(x-1)(x-3)$; máx. local $g(1)=\tfrac43$; mín. local $g(3)=0$; $g(4)=\tfrac43$.
\item[\textbf{5.26}] $2x\,e^{x^4}-e^{x^2}$.
\item[\textbf{5.27}] $\dfrac{\sin x}{1+\cos^2x}$.
\item[\textbf{5.28}] $\tfrac{23}3$.
\item[\textbf{5.29}] Deslocamento 4 m; distância 12 m; $s(4)=6$.
\item[\textbf{5.30}] $f_{\text{méd}}=\tfrac2\pi\approx0{,}637$; $c=\sin^{-1}\tfrac2\pi\approx0{,}690$ ou $c=\pi-0{,}690\approx2{,}451$.
\item[\textbf{5.31}] $\tfrac25(x+2)^{5/2}-\tfrac43(x+2)^{3/2}+C$.
\item[\textbf{5.32}] $-\csc x+C$; $\tfrac14\tan^4x+C$.
\item[\textbf{5.33}] $\tfrac83$.
\item[\textbf{5.34}] $\ln\tfrac{1+e}2\approx0{,}620$.
\item[\textbf{5.35}] $\tfrac13(\ln x)^3+C$; $\ln|\ln x|+C$.
\item[\textbf{5.36}] $\tan^{-1}(e^x)+C$.
\item[\textbf{5.37}] $\tfrac\pi6$; $\tfrac\pi3$.
\item[\textbf{5.38}] $\tfrac12\tan^{-1}\tfrac{x-3}2+C$.
\item[\textbf{5.39}] $0$.
\item[\textbf{5.40}] $\tfrac12\ln(x^2+1)+3\tan^{-1}x+C$.
\item[\textbf{5.41}] $4\tan^{-1}1=\pi$; $M_4\approx3{,}14680$.
\item[\textbf{5.42}] $6000$ L; em $t=0$; $V(t)=1000+200t-2t^2$, parábola crescente até o vértice em $t=50$.
\item[\textbf{5.43}] $\ln(\sqrt2+1)\approx0{,}8814$.
\item[\textbf{5.44}] $10$.
\item[\textbf{5.45}] $f(x)=x^{3/2}$, $a=9$.
\item[\textbf{5.46}] Mín. local em $x=2$; inflexão em $x=2+\sqrt5$.
\item[\textbf{5.47}] $\tfrac13$ (L'Hôpital + TFC).
\item[\textbf{5.48}] $2$; $4\sqrt2$.
\item[\textbf{5.49}] $500(1-e^{-1{,}6})\approx399$; tende a $500$.
\item[\textbf{5.50}] $-\sqrt{1-x^2}+\tfrac13(1-x^2)^{3/2}+C$.
\item[\textbf{5.51}] $\sec^{-1}(e^x)+C$.
\item[\textbf{5.52}] $-\tan^{-1}(\cos x)+C$; $\tfrac12\sec^{-1}\tfrac{|x|}2+C$.
\item[\textbf{5.53}] $\sin^{-1}\tfrac{x-2}2+C$.
\item[\textbf{5.54}] Derive $G(a)=\int_a^{a+T}f$: $G'(a)=f(a+T)-f(a)=0$.
\item[\textbf{5.55}] $56{,}25$ m.
\item[\textbf{5.56}] $\tfrac1{\sqrt2}\approx0{,}707$.
\item[\textbf{5.57}] $g''=f'\ge0$; $g$ acumula valores $f(t)\ge0$.
\item[\textbf{5.58}] $\int_0^1\sqrt x\,dx=\tfrac23$.
\item[\textbf{5.59}] $\int_0^1\tfrac{dx}{1+x^2}=\tfrac\pi4$.
\item[\textbf{5.60}] $\tfrac23$; $2\sqrt2$.
\item[\textbf{5.61}] A integral $I$ satisfaz $I=\int_0^{\pi/2}\tfrac{\cos^n}{\sin^n+\cos^n}$; some as duas: $2I=\tfrac\pi2$.
\item[\textbf{5.62}] $\int_a^{ab}\tfrac{dt}t=\int_1^b\tfrac{du}u$ com $t=au$.
\item[\textbf{5.63}] $\tfrac\pi8\ln2$.
\item[\textbf{5.64}] Nenhuma: derivando, $f=\tfrac{2x}{1+x^2}$, mas em $x=0$ seria preciso $\int_0^1t^2f=0$, e essa integral vale $1-\ln2$.
\item[\textbf{5.65}] Retângulos acima e abaixo de $y=1/x$ em $[1,n+1]$ e $[1,n]$.
\end{itemize}\end{multicols}}

